\section{Analytical inverse kinematics}
\label{sec:analytical-ik}

\subsection{Frames, geometry and the wrist target}
\label{sec:ik-position}
We seek the joint vectors $\qq=(q_1,\ldots,q_6)^\trans$ giving a prescribed pose
of \texttt{link\_6}. Frame $0$ is \texttt{base\_link}; frames $1,\ldots,6$
are the URDF link frames. The $q_i$ are URDF joint displacements from home.

Let $C$ be the wrist centre, the origin of frame~5. The geometry is specified by
\begin{equation}
 \begin{aligned}
 {}^0\vect t_2\big|_{q_1=0}&=\begin{pmatrix}a_1\\0\\d_1\end{pmatrix},
 &{}^2\vect t_3&=\begin{pmatrix}a_2\\d_2\\0\end{pmatrix},\\[3pt]
 {}^3\vect p_C&=\begin{pmatrix}d_4\\0\\d_3\end{pmatrix},
 &Q:={}^{1}R_2\big|_{q_2=0}&=R_x(\alpha).
 \end{aligned}
 \label{eq:ik-geometry}
\end{equation}
Here $d_1$ combines the two vertical base offsets, and $d_4$ combines the
$x$ offsets from joint~3 to joint~4 and from joint~4 to joint~5.
For the supplied robot, $(a_1,a_2,d_1,d_2,d_3,d_4,d_6)
=(0.8,2,1,1.25,0.75,2.5,1.5)\,\mathrm m$ and
$\alpha=\pi/6$. Write $c_i=\cos q_i$, $s_i=\sin q_i$, $c_\alpha=\cos\alpha$ and
$s_\alpha=\sin\alpha$. The matrices $R_x,R_y,R_z$ are right-handed rotations;
$\ex=(1,0,0)^\trans$ and $\ez=(0,0,1)^\trans$ are coordinate unit vectors.
The two-argument function $\operatorname{atan2}(s,c)$ denotes the four-quadrant
inverse tangent, with the sine argument first.

All three wrist axes pass through $C$, so its position depends only on
$q_1,q_2,q_3$. The origin of frame~6 lies a distance $d_6$ from $C$ along
its own $x$ axis. The subscript $d$ denotes a desired quantity, so the target
pose gives the desired wrist position
\begin{equation}
 \boxed{{}^0\vect p_{C,d}={}^0\vect t_{6,d}-d_6\,{}^0R_{6,d}\ex
       =\begin{pmatrix}x_d&y_d&z_d\end{pmatrix}^{\trans}}.
 \label{eq:ik-wrist-target}
\end{equation}

\subsection{Writing the wrist-centre position}
Combining the relative translations and rotations gives us
\begin{equation}
 \begin{aligned}
 {}^0\vect p_C(\qq)={}&\underbrace{d_1\ez}_{{}^0\vect p_O}+\underbrace{R_z(q_1)}_{{}^0R_1}
 \underbrace{\Biggl[
   \underbrace{a_1\ex}_{{}^1\vect t_2-{}^1\vect p_O}
   +\underbrace{Q R_y(q_2)}_{{}^1R_2}
   \underbrace{\Biggl(
     \underbrace{\begin{pmatrix}a_2\\d_2\\0\end{pmatrix}}_{{}^2\vect t_3}
     +\underbrace{R_z(q_3)}_{{}^2R_3}
      \underbrace{\begin{pmatrix}d_4\\0\\d_3\end{pmatrix}}_{{}^3\vect p_C}
   \Biggr)}_{{}^2\vect p_C}
 \Biggr]}_{{}^1\vect p_C-{}^1\vect p_O}.
 \end{aligned}
 \label{eq:ik-position-map}
\end{equation}
This is also the position part of the product of exponentials. We evaluate it
from the inside out. First, joint~3 gives us the shoulder-to-wrist coordinates
in frame~2: ${}^2\vect p_C=(F_1,B,F_2)^\trans$, where
\begin{equation}
 F_1:=a_2+d_4c_3,\qquad B:=d_2+d_4s_3,\qquad F_2:=d_3.
 \label{eq:ik-auxiliary}
\end{equation}
Next we apply joint~2, then the fixed tilt and the shoulder offset:
\begin{equation}
 \begin{aligned}
 \underbrace{\begin{pmatrix}X\\B\\D\end{pmatrix}}_{\text{before the tilt}}
 &=R_y(q_2)\begin{pmatrix}F_1\\B\\F_2\end{pmatrix}
 =\begin{pmatrix}F_1c_2+F_2s_2\\B\\F_2c_2-F_1s_2\end{pmatrix},\\[8pt]
 \underbrace{\begin{pmatrix}A\\Y\\H\end{pmatrix}}_{\text{before joint 1}}
 &=a_1\ex+Q\begin{pmatrix}X\\B\\D\end{pmatrix}
 =\begin{pmatrix}a_1+X\\c_\alpha B-s_\alpha D\\s_\alpha B+c_\alpha D\end{pmatrix}.
 \end{aligned}
 \label{eq:ik-intermediate-vectors}
\end{equation}
$A,Y,H$ are the wrist's two horizontal components and its height relative
to the point on axis~1 at height $d_1$, in axes that rotate with joint~1. Applying the final base rotation gives us
\begin{equation}
 x_d=Ac_1-Ys_1,\qquad y_d=As_1+Yc_1,\qquad z_d=d_1+H.
 \label{eq:ik-components}
\end{equation}

\subsection{Eliminating the first two joint angles}
Let $O$ be the point on axis~1 at shoulder height,
${}^0\vect p_O=(0,0,d_1)^\trans$. We now form three \emph{known target quantities}:
\begin{equation}
 h:=z_d-d_1,\qquad \rho^2:=x_d^2+y_d^2,\qquad
 \Rcal:=\rho^2+h^2=\|{}^0\vect p_{C,d}-{}^0\vect p_O\|^2.
 \label{eq:ik-invariants}
\end{equation}
$h$ is the wrist height above $O$, $\rho$ is wrist distance to axis~1, and $\Rcal$ is
the wrist squared distance to $O$. None changes when joint~1 rotates. For a solution,
\eqref{eq:ik-components} gives $H=h$ and $\rho^2=A^2+Y^2$.

Rotation preserves norm. We apply that to the first relation in \eqref{eq:ik-intermediate-vectors}:
\begin{equation}
 \underbrace{X^2+B^2+D^2}_{\|R_y(q_2)(F_1,B,F_2)^\trans\|^2}
 =\underbrace{F_1^2+B^2+F_2^2}_{\|(F_1,B,F_2)^\trans\|^2}
 \quad\Longrightarrow\quad X^2+D^2=F_1^2+F_2^2.
 \label{eq:ik-rotation-norm}
\end{equation}
The unchanged $y$ component $B$ cancels. Likewise,
$(Y,H)^\trans=\bigl(\begin{smallmatrix}c_\alpha&-s_\alpha\\s_\alpha&c_\alpha\end{smallmatrix}\bigr)(B,D)^\trans$
gives $Y^2+h^2=B^2+D^2$. Hence
\begin{equation}
 \begin{aligned}
 \Rcal&=(a_1+X)^2+B^2+D^2=2a_1X+F_3,\qquad
 h=s_\alpha B+c_\alpha D,\\
 F_3&:=a_1^2+F_1^2+F_2^2+B^2
     =L+2a_2d_4c_3+2d_2d_4s_3,\\
 L&:=a_1^2+a_2^2+d_2^2+d_3^2+d_4^2.
 \end{aligned}
 \label{eq:ik-invariant-equations}
\end{equation}
The expansion of $F_3$ uses $c_3^2+s_3^2=1$: $F_3$ depends only on $q_3$,
while $L$ is constant. With $\Rcal,h$ known and $a_1,c_\alpha$ nonzero,
a trial $q_3$ fixes
\begin{equation}
 X=\frac{\Rcal-F_3}{2a_1},\qquad
 D=\frac{h-s_\alpha B}{c_\alpha}.
 \label{eq:ik-XD}
\end{equation}
These required $X,D$ can come from a real rotation $q_2$ only if their norm
is correct. Substituting into \eqref{eq:ik-rotation-norm} and multiplying by
$4a_1^2$ gives us the \emph{closure condition}:
\begin{equation}
 \boxed{G(q_3):=(\Rcal-F_3)^2+
       \frac{4a_1^2}{c_\alpha^2}(h-s_\alpha B)^2
       -4a_1^2(F_1^2+d_3^2)=0.}
 \label{eq:ik-closure}
\end{equation}
Only $q_3$ remains unknown. For $\alpha=0$, $D=h$ and this becomes the
lecture's untilted closure, with terminal arm length $d_4$.

\subsection{Converting the closure into a quartic}
We use the half-angle substitution
\begin{equation}
 t=\tan(q_3/2),\qquad \Delta=1+t^2,\qquad
 c_3=\frac{1-t^2}{\Delta},\qquad s_3=\frac{2t}{\Delta}.
 \label{eq:ik-half-angle}
\end{equation}
The hat below denotes a numerator after substitution.
Multiplying by the common denominator $\Delta$ gives
\begin{equation}
 \begin{aligned}
 \widehat F_1(t)&:=\Delta F_1=a_2\Delta+d_4(1-t^2),\\
 \widehat B(t)&:=\Delta B=d_2\Delta+2d_4t,\\
 \widehat F_3(t)&:=\Delta F_3=L\Delta+2a_2d_4(1-t^2)+4d_2d_4t.
 \end{aligned}
 \label{eq:ik-numerators}
\end{equation}
Each numerator is quadratic in $t$. Multiplying \eqref{eq:ik-closure} by
$c_\alpha^2\Delta^2$ clears the denominators and yields a polynomial of degree
at most four, because it contains squares of these quadratics:
\begin{equation}
 \boxed{\begin{aligned}
 P(t)={}&c_\alpha^2(\Rcal\Delta-\widehat F_3)^2
       +4a_1^2(h\Delta-s_\alpha\widehat B)^2\\
       &-4a_1^2c_\alpha^2(\widehat F_1^2+d_3^2\Delta^2)=0.
 \end{aligned}}
 \label{eq:ik-quartic}
\end{equation}
Every real root of the polynomial gives a $q_3=2\arctan t$. Since $\Delta>0$ for real $t$,
clearing denominators adds no real roots. We must test $q_3=\pi$ separately in
\eqref{eq:ik-closure}, because its half-angle tangent is undefined. This missing point is included in the maximum of four distinct $q_3$ candidates.

\subsection{Recovering the remaining arm angles}
For each candidate $q_3$, we evaluate \eqref{eq:ik-auxiliary} for $F_1,F_2,B$ and
use \eqref{eq:ik-XD} for $X,D$. Combining their defining equations cancels
$s_2$ in $F_1X+F_2D$ and $c_2$ in $F_2X-F_1D$, giving us
\begin{equation}
 c_2=\frac{F_1X+F_2D}{F_1^2+F_2^2},\qquad
 s_2=\frac{F_2X-F_1D}{F_1^2+F_2^2},\qquad
 q_2=\operatorname{atan2}(s_2,c_2).
 \label{eq:ik-q2}
\end{equation}
The denominator is positive because $F_2=d_3\ne0$, and the closure ensures
$c_2^2+s_2^2=1$. Next we set $A=a_1+X$, $Y=c_\alpha B-s_\alpha D$ and recover
the base angle from the horizontal equations:
\begin{equation}
 q_1=\operatorname{atan2}(\underbrace{Ay_d-Yx_d}_{\rho^2s_1},\;
           \underbrace{Ax_d+Yy_d}_{\rho^2c_1}),\qquad \rho>0.
 \label{eq:ik-q1}
\end{equation}
The common positive factor $\rho^2$ does not change the angle.
If $\rho=0$, the wrist centre lies on axis~1 and $q_1$ is free, in which case we
choose a value in $[-\pi,\pi]$ and solve the wrist for that choice.

\subsection{Recovering the wrist angles}
\label{sec:ik-orientation}
For each arm solution, we undo its orientation using the transpose, then solve
the remaining local $x$--$z$--$x$ wrist rotations:
\begin{equation}
 \begin{aligned}
 {}^0R_3&=R_z(q_1)Q R_y(q_2)R_z(q_3),\\
 {}^3R_{6,d}&=({}^0R_3)^\trans\,{}^0R_{6,d}
             =R_x(q_4)R_z(q_5)R_x(q_6)=:(r_{ij}).
 \end{aligned}
 \label{eq:ik-wrist-rotation}
\end{equation}
Here $r_{ij}$ is entry $(i,j)$ of ${}^3R_{6,d}$. Multiplying the wrist
rotations gives us five useful entries:
\begin{equation}
 r_{11}=c_5,\quad r_{21}=c_4s_5,\quad r_{31}=s_4s_5,\quad
 r_{12}=-s_5c_6,\quad r_{13}=s_5s_6.
 \label{eq:ik-wrist-entries}
\end{equation}
Let $\eta=\sqrt{r_{21}^2+r_{31}^2}=|s_5|$. For $\eta>0$, we choose each sign
$\varepsilon\in\{-1,+1\}$, so that $s_5=\varepsilon\eta$. Then
\begin{equation}
 \boxed{\begin{aligned}
 q_5&=\operatorname{atan2}(\varepsilon\eta,r_{11}),\\
 q_4&=\operatorname{atan2}(\varepsilon r_{31},\varepsilon r_{21}),\qquad
 q_6=\operatorname{atan2}(\varepsilon r_{13},-\varepsilon r_{12}).
 \end{aligned}}
 \label{eq:ik-wrist-solutions}
\end{equation}
We must use the same sign in all three formulas. It selects one of two
wrist configurations with the same orientation. If $\eta=0$, joints~4 and~6
are coaxial, so only a sum or difference is fixed:
\begin{equation}
 \begin{array}{lll}
 r_{11}=1: &q_5=0,&q_4+q_6=\operatorname{atan2}(r_{32},r_{22}),\\[2pt]
 r_{11}=-1:&q_5=\pi,&q_6-q_4=\operatorname{atan2}(r_{32},-r_{22}).
 \end{array}
 \label{eq:ik-wrist-singularity}
\end{equation}
We may set $q_4=0$ and recover $q_6$ from the appropriate relation.

\subsection{Assembling and checking the solutions}
For every real $q_3$ candidate, we recover the arm and both wrist branches.
We represent angles in $[-\pi,\pi]$, enforce joint limits, remove duplicates
modulo $2\pi$, and check position and orientation by forward kinematics.
We use numerical tolerances for roots, singularities and pose errors.
Away from continuous singular families, at most four arm solutions each
have two wrist branches, so there are \textbf{at most eight full-pose solutions}.
